\documentclass[11pt,a4paper]{ltjsarticle}
\usepackage{luatexja-fontspec}
\usepackage{geometry}
\usepackage{booktabs}
\usepackage{longtable}
\usepackage{array}
\usepackage{enumitem}
\usepackage{hyperref}
\usepackage{xcolor}
\usepackage{amsmath}
\usepackage{fancyvrb}
\usepackage{tcolorbox}

\geometry{margin=25mm}

\setmainjfont{Hiragino Mincho ProN}
\setsansjfont{Hiragino Kaku Gothic ProN}
\setmonofont{Menlo}

\hypersetup{
  colorlinks=true,
  linkcolor=blue!60!black,
  urlcolor=blue!60!black,
}

% プロンプト引用用ボックス
\newtcolorbox{promptbox}{
  colback=gray!5,
  colframe=blue!40!black,
  boxrule=0.5pt,
  left=8pt, right=8pt, top=6pt, bottom=6pt,
  fontupper=\small\ttfamily,
}

\title{{\sffamily 進捗報告：変数共有関係を表す数式グラフを用いた\\GNNに基づく文献からの物理モデル構築}}
\author{かずひろ}
\date{2026年4月13日（4月17日改訂）}

\begin{document}
\maketitle


%======================================================================
\section{研究目的}
%======================================================================

科学文献PDFから数式を自動抽出し，自然言語で記述されたプロセス条件（例：「CSTRの定常濃度を求めたい．入力は流量と反応速度定数」）に対して，最適な物理モデル（数式の集合）を自動選択するRetrievalシステムを開発する．数式間の変数共有関係をグラフ構造で表現し，グラフニューラルネットワーク（GNN）ベースの検索手法の有効性を検証する．

%======================================================================
\section{使用データ}
%======================================================================

\subsection{データ概要}

\begin{table}[h]
\centering
\begin{tabular}{ll}
\toprule
項目 & 値 \\
\midrule
総数式数 & 3,244 \\
ソース数 & 71（PDF 39 + ハンドブック 32） \\
総変数ノード & 3,704 \\
グラフエッジ（二部） & 33,638 \\
訓練ケース & 1,107（123コア $\times$ 9バリアント） \\
\bottomrule
\end{tabular}
\end{table}

\textbf{変数の表記揺れへの対処方針}：
異なる文献で同一の物理量が異なる記号で表記されるケース（例：比熱を$c_p$と書く文献と$c_w$と書く文献）が存在する．前回報告時に「変数表記の正規化に取り組む」旨を述べたが，実データを精査したうえで，\textbf{現段階では全面的な正規化を行わず，原典表記を保持する方針}を採る判断をした．以下に実態と判断根拠を示す．

\textbf{表記揺れの実態}（本研究で構築した3,244式の集計）：

\begin{table}[h]
\centering
\small
\begin{tabular}{lll}
\toprule
物理量 & 出現している表記（回数） & 合計 \\
\midrule
比熱 & $c_p$(18), $C_p$(78), $c_w$(2), $Cp$(1) & 99 \\
前指数因子 & $k_0$(23), $k_o$(6), $k0$(2) & 31 \\
密度 & \verb|\rho|(242), $\rho$(7, Unicode), rho(2) & 251 \\
\bottomrule
\end{tabular}
\end{table}

\noindent 同一物理量が最大4種類の記号で分散しており，表記揺れは定量的に存在する．

\textbf{全面正規化を行わない判断の根拠}：

\begin{enumerate}[leftmargin=2em, nosep]
  \item \textbf{誤った変数統合のリスク}：記号$P$は文脈により圧力・密度・電力・確率など多義的に使われ，$Q$は流量・熱量・電荷，$k$は反応速度定数・熱伝導率・ばね定数と，多義の例は多い．文脈を見ずに機械的に統一すると，\textbf{異なる物理量を同一視する誤り}が生じる．
  \item \textbf{評価設定における効果の限定性}：本研究では訓練・評価データを文献単位で分割している．同一文献内では著者が一貫した表記を用いるため，クエリと正解式の変数表記は多くの場合既に一致している．正規化が効くのは同一ドメイン内で複数の文献にまたがる正解式が与えられるケースに限定され，全体への効果は限定的と見込まれる．
  \item \textbf{現時点のMRRが既に高水準}：reranker-10で$0.949\pm0.031$に達しており，表記揺れに起因する誤りが残り誤差（5.1\%）の主要因でない可能性が高い．誤正規化のリスクを負って全面適用するよりも，他の改善手段（hard negative生成，訓練ケース増強等）を優先する方が費用対効果が高いと判断した．
\end{enumerate}

\textbf{今後の方針}：
ルールベースでの明示的な記号統一には上記の多義性問題があるため，将来的には，各変数に対して抽出時に付与している物理的意味の自然言語記述（例：「density of fluid j [kg/m$^3$]」）を埋め込み空間で比較し，\textbf{意味的に近い変数ノードをグラフ上でソフトに接続}する学習的アプローチを検討する．これにより，記号レベルの完全一致を超えて，物理的意味レベルでの変数関連を捉えることを目指す．

\subsection{参考文献リスト（PDF資料）}

主要なPDFソースから抽出した数式（合計2,574式）．

\textbf{注}：同一書籍を複数行に分けて記載しているものは，Gemini APIのページ上限のためにPDFを章単位で分割して抽出した結果であり，同一書籍の異なるチャンクに対応する（\#1と\#2，\#8と\#9）．

\begin{longtable}{clrc}
\toprule
\# & ソース & 式数 & 分野 \\
\midrule
\endhead
1 & Seborg et al., \textit{Process Dynamics and Control}, 3rd Ed. & 308 & プロセス制御 \\
2 & \quad 同上，Chap.~2（別チャンクで分割抽出） & 108 & プロセス制御 \\
3 & Bequette, \textit{Process Control: Designing Processes} & 482 & プロセス制御 \\
  & \textit{and Control Systems for Dynamic Performance} & & \\
4 & Woolf et al., \textit{Chemical Process Dynamics and Controls} & 521 & 化学プロセス動特性 \\
5 & Ingham et al., \textit{Chemical Engineering Dynamics} (Wiley) & 70 & 化学工学動特性 \\
6 & Marlin, \textit{Process Control}, Chapter 3--4 & 161 & プロセス制御 \\
7 & Svrcek et al., \textit{A Real Time Approach to Process Control} & 61 & プロセス制御 \\
8 & Corriou, \textit{Process Control} (Springer, 2018) & 226 & プロセス制御 \\
9 & \quad 同上，Chap.~1/1--6（別チャンクで分割抽出） & 104 & プロセス制御 \\
10 & Aris (1974), 反応系の合理的解析 & 33 & 反応工学 \\
11 & Moser et al.\ (2018), ワイン発酵モデル & 61 & バイオプロセス \\
12 & Liu \& Guo (2013), プレウロムチリン発酵モデル & 15 & 発酵工学 \\
13 & Phisalaphong et al.\ (2006), エタノール発酵最適化 & 15 & 発酵工学 \\
14 & Bhandari et al.\ (2009), バイオプロセス速度論 & 9 & バイオプロセス \\
15 & 燃焼関連論文3本（2023--2024）\textsuperscript{※,\,†} & 79 & 燃焼工学 \\
16 & 高分子反応器論文\textsuperscript{※} & 9 & 高分子工学 \\
17 & 熱交換器論文（動モデル）\textsuperscript{※} & 71 & 伝熱工学 \\
18 & 物質移動・対流論文\textsuperscript{※} & 42 & 輸送現象 \\
19 & 熱分解・電気化学・大豆油速度論文\textsuperscript{※} & 59 & 反応工学・電気化学 \\
20 & applsci-10-00992-v2（\textit{Applied Sciences}, 2020），バイオディーゼル製造プロセスのエステル交換反応モデル\textsuperscript{※} & 42 & 反応工学・バイオ燃料 \\
21 & energies-17-01344（\textit{Energies}, 2024），エステル交換プロセスの動特性モデル\textsuperscript{※} & 23 & 反応工学・バイオ燃料 \\
22 & Chemical Process Dynamics（書誌情報未特定）\textsuperscript{※}，反応速度定数モデル & 23 & 反応工学 \\
23 & 038.pdf（書誌情報未特定）\textsuperscript{※}，熱伝達を伴う複合CSTR装置モデル & 13 & プロセス制御 \\
24 & \textit{Chemical reactor stability research of ethylene polymerization} (2023)\textsuperscript{※} & 11 & 高分子反応工学 \\
25 & MATEC Web of Conferences (CSCC 2019)\textsuperscript{※}，エステル交換反応の速度論 & 9 & 反応工学 \\
26 & vol20no2p7.pdf（書誌情報未特定）\textsuperscript{※}，連続攪拌槽反応器モデル & 8 & プロセス制御 \\
27 & \textit{World Journal of Engineering and Technology} (2015)\textsuperscript{※}，CSTRの速度論 & 6 & プロセス制御 \\
28 & \textit{An experimental and modeling study of ammonia with enriched oxygen content and ammonia/hydrogen laminar flames} (2021)\textsuperscript{※} & 3 & 燃焼工学 \\
29 & Process Control Symbols and Acronyms（補助資料）\textsuperscript{※}，化学反応速度論の記号定義 & 2 & プロセス制御 \\
\bottomrule
\end{longtable}

\noindent\textsuperscript{※}\textbf{注}：\#15--29はSemantic Scholar APIにより自動取得した論文群である．\#15--19は分野ごとにまとめて記載しており，個別の著者名・タイトルは次回報告で記載する．\#20--29は今回個別に記載したが，\#22・\#23・\#26については書誌情報（著者・出版年）が特定できておらず，次回までに精査する．

\noindent\textsuperscript{†}\textbf{燃焼論文が3本である理由}：本研究では各カテゴリについてSemantic Scholar APIで20--30件を検索し，(i) PDFリンクが存在すること，(ii) Open-accessで取得可能であること，(iii) 数式抽出に十分な数学的記述を含むこと，の条件で絞り込んだ．燃焼カテゴリでは上記条件を満たしたのが3本であった（2023--2024年の乱流予混合燃焼・火炎伝播に関する論文群で，合計79式を抽出）．件数自体は事前に決めた値ではなく，open-accessかつ数式含有量の多い論文を採った結果である．他カテゴリ（\#16--20）も同様のフィルタリング結果として件数が決まっている．

\subsection{ハンドブック式（LLM生成，合計670式）}

\textbf{「ハンドブック式」の定義}：
ハンドブック式とは，Perry's Chemical Engineers' Handbook等の工学ハンドブックに広く収録されている基礎方程式（例：Navier--Stokes方程式，Bernoulliの式，Fourier の法則）を指す．これらは特定の論文に帰属しない普遍的な工学知識であり，2.2節のPDFソース（特定の教科書・論文から抽出した式）とは性質が異なる．PDFソースの数式はプロセス制御に偏っているため，ハンドブック式を追加することで分野の網羅性を補強する目的がある．

本研究では電子版ハンドブックを入手できなかったため，Google Gemini 2.5 Proの知識ベースから基礎方程式を生成した．

\textbf{ドメインの選択基準}：
プロセスシステム工学の主要分野を網羅するため，化学工学の標準的なカリキュラムおよびPerry's Handbookの章構成を参考に，Round~1で基幹15分野を選定した．Round~2では，Round~1で未カバーの応用分野（原子力，半導体，航空など）に加え，Round~1の主要分野の応用的な方程式（例：流体力学の基礎式に対する乱流モデルや境界層の式）を追加した．

\textbf{Round 1（15ドメイン，基礎式）}：
流体力学(25式)，伝熱(25)，物質移動(20)，熱力学(25)，反応工学(25)，プロセス制御(25)，電気化学(20)，バイオプロセス(20)，環境工学(20)，振動学(20)，電気工学(20)，高分子工学(20)，燃焼工学(20)，分離プロセス(20)，輸送現象(25)

\textbf{Round 2（16ドメイン，応用式）}：
Round~1の分野ではカバーしきれなかった応用方程式および新規分野を追加．
原子力工学(20)，半導体物理(20)，音響学(20)，光学(20)，地盤工学(20)，航空力学(20)，HVAC(20)，食品工学(20)，製薬工学(20)，構造力学(20)，化学工学基礎(20)，流体力学応用(20)，伝熱応用(20)，反応工学応用(20)，熱力学応用(20)，電気化学応用(20)，プロセス制御応用(20)

\subsection{Google Gemini API 利用に関する確認事項}

本研究では数式抽出および訓練ケース生成にGoogle Gemini API（gemini-2.5-pro）を使用している．利用規約上の確認結果を以下に示す．

\begin{table}[h]
\centering
\begin{tabular}{lll}
\toprule
項目 & 有料枠 (Paid) & 無料枠 (Unpaid) \\
\midrule
出力の所有権 & Googleは所有権を主張しない & 同左 \\
モデル改善への使用 & 使用されない & 使用される可能性あり \\
人間レビュー & なし & 確認される可能性あり \\
\bottomrule
\end{tabular}
\end{table}

\textbf{対応方針}：
\begin{itemize}[nosep]
  \item 研究段階では\textbf{有料枠を使用する}．無料枠ではプロンプトや出力がGoogleのモデル改善に利用される可能性があり，研究内容の流出リスクがある．
  \item \textbf{PDFアップロード時の著作権リスク}：無料枠（データがモデル学習に使用される設定）で著作権のあるPDFをアップロードした場合，著作権者から訴えられるリスクがある．有料枠では入力データがモデル改善に使用されないため，このリスクは軽減される．
  \item Googleの利用規約上，APIの出力結果を研究に使用すること自体は禁止されていない（出力の所有権はユーザーに帰属）．
  \item AI生成コンテンツの利用開示については，各学会・論文誌の規定に従う．
\end{itemize}

\textbf{参考URL}：
\begin{itemize}[nosep]
  \item Gemini API 利用規約: \url{https://ai.google.dev/gemini-api/terms}
  \item Generative AI 追加規約: \url{https://policies.google.com/terms/generative-ai}
  \item データログポリシー: \url{https://ai.google.dev/gemini-api/docs/logs-policy}
\end{itemize}

%======================================================================
\section{情報抽出手法}
%======================================================================

\subsection{PDFから数式の構造化抽出}

\textbf{使用モデル}: Google Gemini 2.5 Pro（Structured Output モード）

PDFファイルをGemini APIにアップロードし，以下のプロンプトで数式を構造化JSONとして抽出する．出力スキーマはPydanticモデルで定義し，\texttt{response\_json\_schema}パラメータで型安全なJSON出力を強制している．

\textbf{注意}：著作権のあるPDFのアップロードに際しては，有料枠（入力がモデル学習に使用されない設定）を使用する．

\textbf{入力}: PDF文書ファイル \quad
\textbf{出力}: 数式のJSON配列（各要素: source\_id, eq\_id, equation, variables, context\_text, domain）

\begin{promptbox}
Extract all numbered equations and equations found in the appendices
from this textbook in LaTeX format. For each equation, identify the
"variable definitions" and the "physical meaning (context)" from the
surrounding text.

Constraints:
- Extract the content EXACTLY as it appears in the literature.
  Do not make unauthorized corrections or modifications.
- Preserve every variable symbol and subscript exactly as printed:
  e.g.\ if the source writes c\_w (specific heat of water),
  do NOT change it to c\_p.
  Never "normalize" or "standardize" notation.

Output a single JSON array. Each element must have these fields:
- source\_id: string (e.g.\ PDF filename)
- eq\_id: string (e.g.\ "eq\_1")
- equation: string (LaTeX form)
- variables: object \{symbol: definition\}
- context\_text: string (physical meaning)
- domain: string (e.g.\ "CSTR Kinetics")
\end{promptbox}

\textbf{Gemini APIの制約と対処}：
\begin{itemize}[nosep]
  \item レート制限: 5 RPM（12秒間隔で呼び出し）
  \item ページ上限: 約1,000ページ．大型テキストブック（3冊，最大1,400ページ）はpypdfで80ページ単位のチャンクに分割し，各チャンクを個別に処理（合計39チャンク）
\end{itemize}

\subsection{ハンドブック式のLLM生成}

\textbf{入力}: ドメイン名，トピックリスト，生成式数 \quad
\textbf{出力}: 数式のJSON配列（抽出と同一スキーマ）

\textbf{ドメイン名・トピックリストの準備方法}：現時点では，対象分野のカリキュラムやハンドブックの目次を参考に手動で設計している．将来的には，対象プロセスの記述から自動的に関連ドメインとトピックを推定するパイプラインの構築が必要であるが，現段階ではデータベース構築のために網羅的に主要分野を準備する方針を取っている．

\begin{promptbox}
You are an expert in \{domain\_label\} and mathematical modeling.
Generate exactly \{n\} fundamental equations used in \{domain\_label\},
in LaTeX format.

Cover these topics: \{topics\}

For each equation provide:
- source\_id: use exactly "\{source\_id\}"
- eq\_id: "eq\_1" through "eq\_\{n\}" (sequential)
- equation: LaTeX form (use standard notation)
- variables: object \{symbol: definition with units\}
- context\_text: brief physical meaning (1--3 sentences)
- domain: "\{label\}" or a more specific sub-domain

IMPORTANT:
- Each equation MUST be distinct (no duplicates).
- Include differential and algebraic forms where applicable.
- Variables dict must include ALL symbols in the equation.
\end{promptbox}

\textbf{トピックリストの例（流体力学）}: Navier--Stokes equations, Bernoulli equation, Hagen--Poiseuille flow, Darcy--Weisbach friction, boundary layer equations (Blasius), drag and lift coefficients, continuity equation, Euler equation, potential flow, vorticity equation, turbulent flow, open channel flow (Manning, Chezy), orifice flow, Torricelli's theorem

\subsection{論文の自動取得}

\textbf{Semantic Scholar API}により15カテゴリのクエリで関連論文を自動検索し，Open-access PDFをダウンロード．

\textbf{検索カテゴリの設計基準}：PDFソース（2.2節）がプロセス制御に偏っているため，不足する分野（燃焼，高分子，伝熱，電気化学等）をカバーするクエリを設計した．各カテゴリで20--30件を検索し，PDFリンクのあるOpen-access論文のみを対象とした．カテゴリ数・検索件数は，API制限と所要時間のバランスから経験的に設定した値であり，最適化は今後の課題である．

\textbf{検索クエリの例}: ``fluid dynamics pipe flow mathematical model equations'', ``electrochemical cell battery model equations'', ``bioprocess fermentation kinetic model mathematical'' 等

\textbf{結果}: 180論文発見 $\to$ 49論文にPDFリンクあり $\to$ 17本ダウンロード成功 $\to$ 308式抽出

\subsection{訓練ケースの生成}

\subsubsection{コアケース生成}

3,244式の数式一覧（851K文字）をドメイン別に300式以下のバッチに分割し，各バッチからGemini 2.5 Proで物理的に意味のあるケースを生成．

\textbf{入力}: 数式一覧（model\_id, domain, variable\_symbolsのリスト）

\textbf{出力}: コアケースのJSON配列（各要素: case\_id, context, input\_variables, output\_variables, correct\_model\_ids）

\begin{promptbox}
You are helping to create training data for a graph neural network
that selects relevant equations for physical process modeling.

CATALOG (\{N\} equations, domains: \{domain\_list\}):
\{catalog\_json\}

TASK: Using ONLY the models listed above, create approximately
\{n\_cases\} realistic and physically meaningful core training cases.

For each core case:
- Pick one or more model\_ids that represent a coherent physical model.
- Define a context string describing the physical scenario.
- Choose input\_variables and output\_variables from variable\_symbols.
- correct\_model\_ids: list forming a correct mathematical model.

CONDITIONS FOR correct\_model\_ids:
(1) Given input\_variables and the equations, it MUST be possible
    to solve for all output\_variables.
(2) The equation set MUST be sufficient and minimal.
\end{promptbox}

\textbf{結果}: 11バッチ中9バッチ成功（2バッチはGemini 503エラーで失敗）$\to$ 123コアケース

\subsubsection{データ拡張}

各コアケースから9バリアントをPythonテンプレートで自動生成．

\begin{table}[h]
\centering
\begin{tabular}{llrl}
\toprule
バリアント & 手法 & 件数 & 目的 \\
\midrule
original & そのまま使用 & 123 & ベースケース \\
context\_paraphrased & テンプレートで文脈を言い換え（3パターン） & 369 & 語彙の多様性向上 \\
swap\_io & 入力と出力をスワップ\textsuperscript{※1} & 123 & 逆問題への対応力 \\
random\_io\_from\_models & 変数からランダムにI/O割当\textsuperscript{※2} & 492 & 変数選択の多様性 \\
\midrule
\textbf{合計} & & \textbf{1,107} & \\
\bottomrule
\end{tabular}
\end{table}

\textsuperscript{※1}\textbf{swap\_ioの制約}：入力変数数と出力変数数が等しいケースのみスワップを適用している．変数数が異なる場合，スワップ後に自由度の不一致が生じ，数学的に解けないモデルとなるため除外している．

\textsuperscript{※2}\textbf{random\_ioの制約}：ランダムI/O割当後に，選択された数式群で出力変数が実際に求解可能かの検証は現時点では実施していない．このため，物理的に解けないケースが含まれる可能性がある．今後，自由度解析（変数数$-$独立方程式数$=0$の検証）を追加し，不正なケースを除外する仕組みの導入を検討する．

\textbf{言い換えテンプレート}：
\begin{enumerate}[nosep]
  \item ``This case describes the following physical model: \{context\}''
  \item ``In this scenario, the underlying process can be summarized as: \{context\}''
  \item ``From the perspective of process modeling, this case focuses on: \{context\}''
\end{enumerate}

%======================================================================
\section{検索手法の設計}
%======================================================================

本節では，評価結果（第5節）で比較する3手法（baseline, reranker-7, reranker-10）の構成を説明する．3手法はいずれも\textbf{2段階アーキテクチャ}（候補抽出 $\to$ ランキング）に基づく．

\begin{table}[h]
\centering
\begin{tabular}{llll}
\toprule
結果表での名称 & Stage 1（候補抽出） & Stage 2（ランキング） & 学習 \\
\midrule
baseline & TF-IDF+Jaccard（Top 50） & なし（Stage 1のスコアで確定） & なし \\
reranker-7 & 同上 & MLP（7次元特徴量） & あり \\
reranker-10 & 同上 & MLP（10次元特徴量=7+グラフ3） & あり \\
\bottomrule
\end{tabular}
\end{table}

\subsection{用語の説明}

本研究で用いる基本的な手法・指標を説明する．

\begin{description}[style=nextline, leftmargin=2em]
  \item[TF-IDF（Term Frequency--Inverse Document Frequency）]
    文書中の各単語の重要度を数値化する手法．\textbf{Term}（単語）は，数式のLaTeXテキスト・物理的意味・ドメイン名を連結した文字列をトークン化したもの．\textbf{Document}（文書）は，各数式1つずつに対応する上記の連結文字列．頻出する一般的な単語は重み小，特定の数式にしか現れない専門用語は重み大となり，クエリと数式のテキスト的な類似度を計算する．
  \item[SVD（Singular Value Decomposition, 特異値分解）]
    高次元のTF-IDFベクトル（数千次元）を低次元（256次元）に圧縮し，意味的に近い単語パターンを統合する次元削減手法．
  \item[Jaccard係数]
    2つの集合$A, B$の類似度を$|A \cap B| / |A \cup B|$で計算する指標．本研究では，クエリが要求する入出力変数の集合と，候補数式が含む変数の集合との重複度を測定するために使用する．
  \item[L2正規化]
    ベクトルの長さ（L2ノルム）を1に正規化する処理．これにより，ベクトルの「方向」のみで類似度を比較でき，大きさの違いの影響を除去する．
  \item[cosine類似度]
    2つのベクトル間の角度に基づく類似度．$\cos(\mathbf{a}, \mathbf{b}) = \mathbf{a} \cdot \mathbf{b} / (|\mathbf{a}||\mathbf{b}|)$．L2正規化済みベクトルでは内積と一致する．
\end{description}

\subsection{数式グラフの構築}

数式と変数の関係を\textbf{二部グラフ}として表現する．二部グラフとは，ノードが2種類（数式ノードと変数ノード）に分かれ，異なる種類のノード間にのみエッジが存在するグラフである．

\begin{table}[h]
\centering
\begin{tabular}{ll}
\toprule
項目 & 値 \\
\midrule
数式ノード & 3,244 \\
変数ノード & 3,704 \\
二部エッジ & 33,638（双方向） \\
ノード特徴量 & 768次元（CodeBERT [CLS] 埋め込み） \\
\bottomrule
\end{tabular}
\end{table}

\textbf{ノード特徴量の計算}：
\begin{itemize}[nosep]
  \item \textbf{数式ノード}: ``Equation: \{LaTeX\} / Context: \{context\_text\}'' をCodeBERT（microsoft/codebert-base）に入力し，[CLS]トークンの768次元埋め込みを使用
  \item \textbf{変数ノード}: その変数を含む全数式のCodeBERT埋め込みの平均
\end{itemize}

\textbf{エッジの定義}: 数式$i$が変数$j$を含む場合，$(i,j)$と$(j,i)$の双方向エッジを張る．

\textbf{グラフの役割}：このグラフは，手法reranker-10でのみ使用される．具体的には，4.5節のグラフ特徴量（2ホップ到達率・近傍重複度・ブリッジ比率）の計算に用いる．2つの数式が同じ変数を共有していれば，変数ノードを介した\textbf{2ホップのパス}（数式$\to$変数$\to$数式）で間接的につながる．この間接的な関連を特徴量として利用する．

\subsection{Stage 1: 候補抽出（全手法共通）}

全3手法で共通の候補抽出を行う．各クエリについて以下のスコアで全3,244式をランキングし，上位50件を候補として抽出する．
\begin{equation}
  \mathrm{score} = 0.7 \times \cos_{\mathrm{TF\text{-}IDF}}(\text{クエリ}, \text{数式}) + 0.3 \times \mathrm{Jaccard}(\mathit{IO}_{\mathrm{query}}, V_{\mathrm{eq}})
\end{equation}
ここで，$\cos_{\mathrm{TF\text{-}IDF}}$はクエリのcontext文と数式のテキスト表現のTF-IDF cosine類似度，$\mathrm{Jaccard}(\mathit{IO}_{\mathrm{query}}, V_{\mathrm{eq}})$はクエリの入出力変数集合$\mathit{IO}_{\mathrm{query}}$と数式の変数集合$V_{\mathrm{eq}}$のJaccard係数である．重み(0.7, 0.3)は予備実験で決定した．

\textbf{baseline手法}は，このStage 1のスコアをそのまま最終ランキングとして使用する（Stage 2なし，学習なし）．

\subsection{Stage 2: MLPリランカー（reranker-7 / reranker-10）}

Stage 1で得られた候補50件に対し，MLPリランカーが特徴量ベクトルに基づきスコアを付け直す．

\textbf{入力}：候補数式1件ごとに計算された特徴量ベクトル（reranker-7は7次元，reranker-10は10次元）

\textbf{出力}：スコア（0--1のスカラー値，高いほど正解の可能性が高い）

\textbf{MLPの構造}：
\[
  \text{特徴量}(n\text{次元}) \xrightarrow{\text{Linear}(64)} \text{ReLU} \to \text{Dropout}(0.1) \xrightarrow{\text{Linear}(64)} \text{ReLU} \to \text{Dropout}(0.1) \xrightarrow{\text{Linear}(1)} \sigma
\]
ここで$\sigma$はシグモイド関数（出力を0--1に変換）．

\textbf{訓練データにおける正例・負例}：
\begin{itemize}[nosep]
  \item \textbf{正例}：各訓練ケースの\texttt{correct\_model\_ids}に含まれる数式．3.4節で生成した訓練ケース（1,107件）から取得する．1ケースあたり正解数式は1--数個．
  \item \textbf{負例}：Stage 1のTop 50候補のうち，正解でない数式から選択．正例1件に対し最大8件の負例をペアリングする．$\min(8, \text{利用可能数})$としているのは，候補50件中の正解数式が多いケースでは利用可能な負例数が50$-$正解数に制限されるため．実際には正解は通常1--3件であり，負例は47--49件中から8件を選ぶため，少数になる問題は生じない．
\end{itemize}

\begin{table}[h]
\centering
\begin{tabular}{ll}
\toprule
ハイパーパラメータ & 値 \\
\midrule
隠れ層 & 64次元 $\times$ 2層 \\
損失関数 & マージンランキング損失（margin $= 0.1$） \\
負例数 & $\min(8, \text{利用可能数})$/正例 \\
バッチサイズ & 16 \\
エポック数 & 15 \\
最適化 & AdamW (lr $= 10^{-3}$, weight\_decay $= 10^{-4}$) \\
\bottomrule
\end{tabular}
\end{table}

\subsection{特徴量の設計}

\subsubsection{ベース7特徴量（reranker-7で使用）}

\begin{table}[h]
\centering
\begin{tabular}{clp{8cm}}
\toprule
\# & 特徴量名 & 計算方法 \\
\midrule
0 & テキスト類似度 & TF-IDF cosine（クエリcontext, 式のcontext+equation+domain） \\
1 & I/O変数Jaccard & $\mathrm{Jaccard}(\text{クエリ入出力変数}, \text{式の変数集合})$ \\
2 & SVD埋め込み類似度 & TF-IDF $\to$ SVD(256次元)空間でのcosine類似度 \\
3 & 入力変数カバー率 & $|\mathit{input\_vars} \cap \mathit{eq\_vars}| \;/\; |\mathit{input\_vars}|$ \\
4 & 出力変数カバー率 & $|\mathit{output\_vars} \cap \mathit{eq\_vars}| \;/\; |\mathit{output\_vars}|$ \\
5 & 式変数中のクエリ割合 & $|\mathit{IO\_set} \cap \mathit{eq\_vars}| \;/\; |\mathit{eq\_vars}|$ \\
6 & ドメイン一致 & クエリcontextと式domainの部分文字列一致（0/1） \\
\bottomrule
\end{tabular}
\end{table}

\subsubsection{グラフ特徴量（reranker-10で追加される3次元）}

以下の3特徴量は，4.2節の二部グラフを用いて計算する．\textbf{核心：これら3特徴量はすべてクエリに条件付き}（同じ候補式でもクエリの入出力変数が異なれば異なる値をとる）．

\begin{table}[h]
\centering
\begin{tabular}{clp{5.5cm}l}
\toprule
\# & 特徴量名 & 計算方法 & 捉える情報 \\
\midrule
7 & 2ホップ到達率 & クエリI/O変数のうち，変数共有グラフ上で候補式に2ホップ以内で到達可能な割合 & 間接的変数関連 \\
8 & グラフ近傍重複度 & Jaccard（候補式の近傍式集合, クエリ変数を含む式集合） & 構造的近接性 \\
9 & 間接変数ブリッジ比率 & 候補式の変数のうち，クエリ変数には含まれないが，クエリ変数を持つ他式にも出現する変数の割合 & 隠れた接続 \\
\bottomrule
\end{tabular}
\end{table}

\subsection{予備実験：直接GNN学習によるアプローチ}

本研究ではリランカー（4.4節）の他に，GNNで数式の埋め込みを直接学習する2種類のアプローチも試行した．いずれもベースラインを超えられなかったため，主要な比較対象からは除外している．

\textbf{(A) GCN Retriever}：
二部グラフ上で2層のGCN (Graph Convolutional Network) を適用し，数式ノードの埋め込み（768次元）を学習する．入力はCodeBERTによる初期埋め込み，出力はGCNで更新された数式埋め込みである．クエリ埋め込み（TF-IDF $\to$ SVD $\to$ MLP で768次元に変換）との cosine類似度でランキングする．損失関数はSampled InfoNCE（正例1件に対し負例64件）．

\textbf{(B) 残差GCN SVD洗練}：
TF-IDF+SVDによるベースライン埋め込み（256次元，L2正規化済み）を出発点とし，2層GCNで残差的に洗練する．出力は$\mathbf{E}_{\mathrm{refined}} = \mathrm{L2norm}(\mathbf{x} + \alpha \cdot \mathbf{h})$（$\alpha$は学習可能パラメータ，0で初期化）．$\alpha$が小さいうちはベースラインとほぼ同じ挙動をし，学習が進むにつれGCNの修正を徐々に取り込む設計．

\textbf{これらが失敗した原因と教訓}：両手法とも，クエリに依存しない固定的な数式埋め込みを学習するアプローチであった．しかし数式retrievalでは，「どの変数を介してグラフ上で関連しているか」がクエリの入出力変数ごとに異なる．この知見が，クエリ条件付きグラフ特徴量（reranker-10の\#7--9）の設計につながった．

%======================================================================
\section{評価方法と結果}
%======================================================================

\subsection{評価設定}

\begin{table}[h]
\centering
\begin{tabular}{lp{10cm}}
\toprule
項目 & 設定 \\
\midrule
比較手法 & baseline（固定重み），reranker-7（7特徴量MLP），reranker-10（+グラフ3特徴量）．各手法の詳細は4.3--4.5節を参照 \\
評価指標 & MRR (Mean Reciprocal Rank), Recall@$K$ ($K=1,3,5,10$) \\
乱数シード & 5種類（42, 123, 456, 789, 1024） \\
データ分割 & source\_id（論文/ソース）単位で train 60\% / val 20\% / test 20\% \\
統計量 & 5 seed 平均 $\pm$ 標準偏差，Bootstrap 95\%信頼区間（2,000反復） \\
\bottomrule
\end{tabular}
\end{table}

\textbf{評価指標の説明}：
\begin{itemize}[nosep]
  \item \textbf{MRR (Mean Reciprocal Rank)}: 各クエリについて，正解数式が検索結果の何位に現れるかの逆数の平均．MRR$=1.0$は全クエリで正解が1位，MRR$=0.5$は平均2位．
  \item \textbf{Recall@$K$}: 上位$K$件に正解が含まれるクエリの割合．
  \item \textbf{source\_id単位分割}: 同一論文の数式が訓練とテストに分かれることを防ぎ，データリークを防止．
\end{itemize}

\subsection{主要結果：データ規模によるスケーリング効果}

\subsubsection{全体MRR（5 seed 平均 $\pm$ 標準偏差）}

\begin{table}[h]
\centering
\begin{tabular}{rccc}
\toprule
式数 & baseline & reranker-7 & reranker-10 (graph) \\
\midrule
975 & $0.875 \pm 0.088$ & $0.885 \pm 0.119$ & $0.879 \pm 0.123$ \\
1,613 & $0.900 \pm 0.069$ & $0.942 \pm 0.052$ & $0.948 \pm 0.046$ \\
\textbf{3,244} & $0.924 \pm 0.040$ & $0.937 \pm 0.018$ & $\mathbf{0.949 \pm 0.031}$ \\
\bottomrule
\end{tabular}
\end{table}

\noindent 3,244式において，reranker-10（グラフ特徴量あり）が全手法中最高のMRRを達成した．

\subsubsection{95\%信頼区間（3,244式）}

\begin{table}[h]
\centering
\begin{tabular}{lcc}
\toprule
手法 & MRR & 95\% CI \\
\midrule
baseline & $0.924 \pm 0.040$ & $[0.894,\; 0.954]$ \\
reranker-7 & $0.937 \pm 0.018$ & $[0.922,\; 0.951]$ \\
\textbf{reranker-10} & $\mathbf{0.949 \pm 0.031}$ & $\mathbf{[0.924,\; 0.973]}$ \\
\bottomrule
\end{tabular}
\end{table}

\subsubsection{Originalケース（LLM合成でないケース）のみのMRR}

\begin{table}[h]
\centering
\begin{tabular}{rccc}
\toprule
式数 & baseline & reranker-7 & reranker-10 (graph) \\
\midrule
975 & 0.839 & 0.860 & 0.842 \\
1,613 & 0.885 & 0.936 & 0.935 \\
\textbf{3,244} & 0.921 & 0.942 & \textbf{0.952} \\
\bottomrule
\end{tabular}
\end{table}

\subsubsection{グラフ特徴量の効果（reranker-10 $-$ reranker-7）}

\begin{table}[h]
\centering
\begin{tabular}{rrl}
\toprule
式数 & $\Delta$MRR & 解釈 \\
\midrule
975 & $-0.006$ & グラフが害 \\
1,613 & $+0.006$ & ほぼ同等 \\
\textbf{3,244} & $\mathbf{+0.012}$ & \textbf{グラフが有効} \\
\bottomrule
\end{tabular}
\end{table}

\subsubsection{バリアント別MRR（3,244式）}

\begin{table}[h]
\centering
\begin{tabular}{lccc}
\toprule
バリアント & baseline & reranker-7 & reranker-10 \\
\midrule
original & 0.921 & 0.942 & \textbf{0.952} \\
context\_paraphrased & 0.927 & 0.939 & \textbf{0.954} \\
swap\_io & 0.928 & 0.943 & \textbf{0.954} \\
random\_io\_from\_models & 0.923 & 0.933 & \textbf{0.944} \\
\bottomrule
\end{tabular}
\end{table}

\subsubsection{Seed別の詳細（3,244式，baseline）}

\begin{table}[h]
\centering
\begin{tabular}{rrrrr}
\toprule
seed & テスト件数 & MRR & Recall@1 & Recall@3 \\
\midrule
42 & 180 & 0.879 & 0.778 & 1.000 \\
123 & 171 & 0.884 & 0.789 & 1.000 \\
456 & 252 & 0.946 & 0.893 & 1.000 \\
789 & 261 & 0.958 & 0.927 & 1.000 \\
1024 & 99 & 0.955 & 0.909 & 1.000 \\
\bottomrule
\end{tabular}
\end{table}

\noindent\textbf{注}: 全手法・全seedで Recall@3 $= 1.000$（上位3件に正解が必ず含まれる）．

%======================================================================
\section{考察}
%======================================================================

\subsection{データ規模がグラフ構造の価値を決定する}

本研究の最も重要な発見は，グラフ特徴量の有効性がデータ規模に明確に依存するスケーリング則である．

\begin{itemize}
  \item \textbf{975式（スパースグラフ）}: 変数ノードの次数が低く，2ホップ到達率がほぼ0/1に二値化．グラフ特徴量はリランカーにノイズを注入し，性能を悪化させた（$\Delta\mathrm{MRR} = -0.006$）．
  \item \textbf{1,613式（遷移領域）}: グラフがやや密になり，グラフ特徴量が微小ながら正の寄与を示すようになった（$\Delta\mathrm{MRR} = +0.006$）．
  \item \textbf{3,244式（密グラフ）}: 変数ノードの次数が十分に増加し，2ホップ到達率・近傍重複度が連続的な値をとるようになった．グラフ特徴量が候補間の有意な差別化を可能にし，全手法中最高性能を達成（$\Delta\mathrm{MRR} = +0.012$）．
\end{itemize}

この結果は，グラフ構造特徴量には一定のデータ密度の閾値が存在し，約1,600式/2,000変数付近に転換点があることを示唆する．

\subsection{クエリ条件付きグラフ特徴量の有効性}

初期に試行したクエリに依存しないGNN埋め込み（4.6節の(A) GCN Retriever，(B) 残差GCN SVD洗練）はいずれもベースラインを超えられなかった．一方，クエリの入出力変数に条件付けしたグラフ特徴量（reranker-10）は，十分なデータ規模のもとで有効に機能した．

これは，数式retrievalにおいて「どの変数を介してグラフ上で関連しているか」がクエリごとに異なるため，固定的な埋め込みでは捉えきれない情報があることを示唆する．

\subsection{語彙的手法の頑健性とグラフの補完的役割}

TF-IDF + Jaccardベースラインは3,244式でMRR 0.924を達成しており，極めて強い．学習型手法のマージンは$+1.3\%$（reranker-7）から$+2.5\%$（reranker-10）である．

グラフ特徴量は語彙的手法を置き換えるのではなく補完するものとして機能しており，2段階アーキテクチャ（語彙的候補抽出 $\to$ グラフベースリランキング）の妥当性が示された．

\subsection{分散の縮小}

データ拡大により，全手法で5 seed間の分散が大幅に縮小した（baseline: $0.088 \to 0.040$，reranker-7: $0.119 \to 0.018$）．これは評価の統計的信頼性の向上を意味し，データの多様性がモデルの安定性にも寄与することを示す．

%======================================================================
\section{今後の計画}
%======================================================================

\begin{table}[h]
\centering
\begin{tabular}{llp{8cm}}
\toprule
優先度 & 項目 & 目的 \\
\midrule
高 & 論文執筆 & GNNを用いた物理モデル自動構築を主テーマに \\
中 & 訓練ケース増強（$123 \to 200+$コア） & API失敗分2バッチの再実行 \\
低 & GNN学習特徴量の再テスト & GCN-refined SVDを3,244式で再評価 \\
低 & 5,000式超への拡大 & スケーリング曲線の全体像解明，飽和点の同定 \\
\bottomrule
\end{tabular}
\end{table}

\end{document}
